\scanpage{148}\mspage{3}
d'avoir rectifié dans mon esprit une erreur grossière, dû à ma grande ignorance de la top. algébrique et de l'homotopie -- j'étais en effet dans le sentiment que les $H$-espaces satisfaisant des conditions d'associativité et commutativité assez strictes (disons équivalent à des $\unsure{{}'}\!\Omega^{i}X$ avec $i$ arbitrairement grand) correspondaient à des groupes topologiques \unsure{commutatifs|commutatif} (en m'inspirant de quelques analogies \dots). Donc entièrement d'accord avec tes observations p.~5.\unsure{.}

Par contre, je suis toujours intrigué par la question suivante : n'y a-t-il pas un analogue du tapis de Dold-Puppe pour les groupes ($1/2$) simpliciaux (pas néc. commutatifs) et quelle forme prend-il ? À vrai dire, je considère le yoga
\diagram{\begin{tikzcd}[ampersand replacement=\&]
\text{ens. }1/2\text{ simplicial}\arrow[leftrightarrow,r] \& \infty\text{-}\deleted{\unsure{\text{catégorie}}}\text{groupoïde}
\end{tikzcd}}
comme étant essentiellement la version ``set-theoretic'' ultime de Dold-Puppe, qui s'en déduirait en explicitant seulement que les groupes abéliens dans $(\infty\text{-cat})$ ne ``sont autres'' que les complexes de chaines dans $(\mathrm{Ab})$. Il faudrait donc déterminer quels sont les groupes de $(\infty\text{-cat})$ : Je peux te dire ce qu'il en est dans $(1\text{-cat})$ -- ce sera \unsure{dévelop|divulg}%
% archive sheet number: 148
